\paragraph{Voting on predictive loss.}
On voting rows held out from tree fitting, let $y$ contain the scalar outcomes
for regression or the one-hot outcomes for classification. Let $f_T$ collect
the stopped forest's predictions with the inputs in $T$ available. A
classification forest averages probability vectors at its leaves and stopped
nodes. Then squared Euclidean loss, divided by the number of rows, gives the
multiclass Brier score~\citep{gneiting2007proper}. For regression it gives
mean squared error. We score each prediction before averaging over masks.
Scoring the average prediction would omit the variation between masks.

Write $z=y-f_\varnothing$ and $f_T-f_\varnothing=\sum_{H\subseteq T}a_H$.
For a set $U$ of input columns define
\[
 w_U=2\langle z,a_U\rangle-
       \sum_{G\cup H=U}\langle a_G,a_H\rangle ,
\]
where a missing $a_U$ is zero and pairs in the sum are ordered. For a
deployed support $S$, let $E_s(S)$ be the expected excess squared loss
$\|y-f_{S\cap T}\|^2-\|y-f_T\|^2$ under a uniform mask of size $s$.

\begin{proposition}
For each $s=1,\ldots,p$ the excess loss per available fraction is
\[
 V_s(S)=\frac{p}{s}E_s(S)=
 \sum_{U\not\subseteq S}
   \frac{\binom{s-1}{|U|-1}}{\binom{p-1}{|U|-1}}\,w_U .
\]
Every size votes to accept $S$ when
$\min_{1\le s\le p}\{\epsilon\|z\|^2-V_s(S)\}\ge0$.
For a common mask law independent of the row and its label, this is equivalent
to bounding expected excess loss under every exchangeable availability law by
$\epsilon\,\mathbb E[|T|]\,\|z\|^2/p$.
\end{proposition}
\begin{proof}
Expanding the two squares cancels $\|z\|^2$. A linear term disappears when
its required set is omitted from $S$. A pair disappears when $G\cup H$
is omitted. In each case the term enters the comparison only when its entire
required set $U$ is present in $T$, an event with probability
$\binom{s}{|U|}/\binom{p}{|U|}$. Multiplying by $p/s$ gives the displayed
coefficient. Under an exchangeable law, the conditional mask is uniform at
each size of positive mass. Average the conditional bounds over sizes, using
$E_0(S)=0$. Conversely, the law concentrated on any one size tests that size.
\end{proof}

To price the deletion of $j$, add the coefficients of all still-live unions
$U\subseteq S$ containing $j$, grouped by $|U|$. The resulting increment in
$V_s$ uses the same binomial factors. When $j$ is removed, each such union
becomes inactive once and its weight is subtracted from the candidate
coefficients of its remaining members. If $D$ bounds path-group size, the
coefficient updates along a whole deletion path cost
$O(\sum_{U:w_U\ne0}|U|)$ after the unions are formed, with $|U|\le2D$.
Keeping $m$ linear terms and $e$ nonzero Gram entries separately gives the
bound $O(D(m+e))$. Forming the Gram matrix and testing candidates at the
requested sizes incur separate costs.
